\section{Related Works}

\paragraph{Weak Supervision} deep learning approach has achieved remarkable success in several domains beyond NLP \cite{Zhang2022ASupervision}. However, the main bottleneck is collecting massively annotated data. Weak supervision replaces ground-truth annotation with automatic annotation based on heuristic rules, gazetteers, or constraints linguistic rules to address this issue. Some techniques called \textit{distant supervision} exploit semantic links from knowledge bases or ontologies \cite{Lison2021Skweak:NLP}. \cite{Karamanolakis2021Self-TrainingSupervision} proposes an iterative self-training method to combine classic weak supervision and inference of the learning model to extract entities not covered by the initial heuristic rules. In the clinical domain, weak supervision has already been used for specific use cases \cite{Cusick2021UsingIdeation,Fries2021Ontology-drivenRecords,Wang2019ARepresentation}.
\paragraph{Clinical Language Models} In a domain-specific context such as clinical, some specific terms are underrepresented or absent in the general domain. As a result, the clinical NLP community has pretrained Language Models (LMs) \cite{Alsentzer2019PubliclyEmbeddings, Lee2020BioBERT:Mining,Alsentzer2019PubliclyEmbeddings} over domain-specific corpora (i.e. MIMIC-III \cite{Johnson2016MIMIC-IIIDatabase}, Pubmed abstracts extraction). These models could be trained from scratch or from checkpoint to specialize a domain-agnostic model \cite{Gururangan2020DontTasks}.

Though, the performance gains are marginal compared to the general language model. The structure and the abbreviated text present in clinical notes hurt performance.
Instead of pretraining a specialized clinical model, machine learning practitioners can fine-tune agnostic-domain LLMs such as the GPT family of models or T5 on the clinical task. Fine-tuned general-purpose models have proven effective in clinical question-answering, protected health information de-identification, and relation extraction \cite{Lehman2023DoModels}. But this approach requires an important infrastructure and a regular re-finetuning if the data distribution of the EHR shifts. Nevertheless, some LLMs have been trained from scratch over clinical domain-specific notes such as GatorTron \cite{Yang2022ARecords}, BioGPT \cite{Luo2022BioGPT:Mining} or ClinicalT5 \cite{Lu2022ClinicalT5:Text} who achieved promising performance on several tasks.
Additionally, in-context learning with agnostic LLMs such as \instruct \cite{Ouyang2022TrainingFeedback} where no weight is modified shows good results \cite{Agrawal2022LargeExtractors, Brown2020LanguageLearners} and outperforms specialized smaller models on several clinical tasks.
\paragraph{Prompt-based Method}
Prompt-based learning for generative language model treats a downstream task as a language modelling problem where a language model predicts the next tokens of the instruction given a textual prompt in input \cite{Sainz2021LabelExtraction}.

In this paradigm, instead of fine-tuning a model to a downstream task (\textit{"pre-train, fine-tune and predict"}), we want to manipulate the behaviour of a pretrained LM using an appropriate prompt to give the desired output (\textit{"pre-train, prompt and predict"}). The significant advantage of this method is that a pretrained language model trained in an unsupervised fashion could be used for many tasks \cite{Liu2023Pre-trainProcessing}.
This way, \textit{prompt engineering} has been developed to explore the most suitable prompt method applied to a LM to solve a given task.

Among these methods, one of the most popular methods for information extraction, question-answering or sentiment analysis is \textit{in-context learning}. In the clinical domain, some works exist on information retrieval and question-answering. The prompt contains three components: the template or the format of the examples, the set of examples and the ordering of prompts such as present in the figure \ref{fig:prompt}. The aim is to provide some training examples in the prompt before the test example. However, the chosen examples and their ordering and the format could impact performance \cite{Zhao2021CalibrateModelsb}; these three components must be tuned to optimize performance.
% zero-shot -> CoT + Knowledge -> which tasks ?

In another way, we can cite works around the chain of thoughts (CoT). This encourages the LLM to explain its reasoning to get more accurate results, especially in mathematical and logical reasoning \cite{Wei2022ChainModels, Cobbe2021TrainingProblems}. This technique could be used with reason edited manually in the prompt or with two separate prompts where the first involves a reasoning task, then concatenated with the second prompt involving the main tasks.

Other techniques involve generated knowledge similar to CoT. Instead of reasoning, the first prompt generates potentially useful information associated with the tasks concatenated in the final prompt \cite{Liu2022GeneratedReasoning}.


\section{Other Papers}


\subsection{KnowledgeSG https://aclanthology.org/2024.emnlp-main.438}
Background of the study:
The paper addresses the challenge of preserving privacy while fine-tuning large language models (LLMs) on sensitive data. Previous solutions, such as using synthetic data for substitution, struggle to simultaneously improve performance and preserve privacy.

Research objectives and hypotheses:
The paper proposes a novel client-server framework called "KnowledgeSG" to enhance synthetic data quality and improve model performance while ensuring privacy.

Methodology:
On the client side, the framework fine-tunes the local model using differentially private stochastic gradient descent (DP-SGD) to learn local knowledge from private data. On the server side, the framework generates raw synthetic instructions, filters them using a professional model, and then generates accurate responses. The optimized model is then transmitted back to the client.

Results and findings:
The results demonstrate the effectiveness of the proposed framework on both privacy and performance benchmarks in the medical and financial domains. Specifically, KnowledgeSG achieves a relative improvement of 120.39\% over the Non-Private training method in the medical free-form evaluation, even surpassing the professional model.

Discussion and interpretation:
The paper discusses the limitations of existing solutions, such as the trade-off between privacy risk and model performance, and how KnowledgeSG addresses these issues by leveraging knowledge distillation from the professional server-side model.

Contributions to the field:
The paper proposes a novel client-server framework that enhances synthetic data quality and improves model performance while maintaining strict privacy protection.

Achievements and significance:
KnowledgeSG achieves superior performance results while preventing memorization or leakage of the private dataset. The framework significantly reduces the reconstruction rate from 97.13 to 0.87, demonstrating its strong privacy-preserving capabilities.

Limitations and future work:
The paper acknowledges that the effectiveness of KnowledgeSG on general tasks remains to be fully explored. It also mentions that the framework involves more communication and computation cost than Non-Private fine-tuning, which the authors believe is justified by the significant reduction in memorization concerns and performance improvements. Future work includes conducting experiments on more general tasks and optimizing the deployment of KnowledgeSG to make it more compatible and lightweight. 